\documentclass[11pt,letterpaper]{article}
\usepackage[margin=0.85in]{geometry}
\usepackage{lmodern}
\usepackage{hyperref}
\hypersetup{colorlinks=false, urlcolor=black, pdfborder=0 0 0}
\setlength{\parindent}{0pt}
\setlength{\parskip}{0.45em}
\raggedright
\pagestyle{empty}

\begin{document}

\noindent
{\bfseries Ke-Cheng Hsu (Clio)}\\
Ph.D.\ Candidate, Department of Economics, HKUST\\
(852) 6840\,9753 \quad $|$ \quad \href{mailto:khsu@connect.ust.hk}{khsu@connect.ust.hk}\\
\href{https://cliohsu1997.github.io/}{cliohsu1997.github.io} \quad $|$ \quad \href{https://github.com/cliohsu1997}{github.com/cliohsu1997}

\vspace{0.7em}
\noindent
October 21, 2026

\vspace{0.35em}
\noindent
Peixuan Li\\
Associate\\
Econic Partners --- Beijing\\
\href{mailto:pli@econic.com}{pli@econic.com}

\vspace{0.2em}
Dear Peixuan,

I am writing to express my interest in a Ph.D.\ Associate role with Econic Partners in Beijing. Muyang mentioned that the Beijing office is recruiting Ph.D.\ Associates. I am a Ph.D.\ candidate in Economics at HKUST, working in empirical industrial organization and digital economics, and I am a native Mandarin speaker and professionally fluent in English.

My research uses large proprietary datasets to study welfare and run counterfactuals. In ``Too Much Unsold Stock?,'' joint with Kei Kawai (Senior Consultant, Compass Lexecon) and Kohei Kawaguchi (Advisor, UTokyo Economic Consulting Inc.), we ask whether resale price maintenance can be pro-competitive---by encouraging sellers to carry inventory---and how much of the resulting unsold stock is an environmental cost of price commitment. Estimating a dynamic oligopoly model on monthly bookstore data, we find that the classical static model overstates marginal cost by about 20\% when dynamic demand uncertainty and inventory decisions are ignored. Working on this project taught me to let the data speak first---cleaning and transforming raw records, then using descriptive statistics, regression, and clear figures to establish the key facts before turning to structural models.

My job market paper studies return policy as a trade-off between purchase and return efficiency: a stricter return policy curbs inefficient ex-post returns, but the resulting lower return rates make higher prices more attractive for sellers and can exclude ex-ante efficient purchases. Automatic provision of return-shipping insurance performs better in high-margin categories by expanding valuable purchases, while selective adoption outperforms automatic provision by 15--30 percentage points of first-best welfare in thinner-margin categories by curbing excessive returns. A partial return-penalty discount (a discounted price plus a 10--15\% restocking fee) outperforms both outside dresses, attaining 75--78\% of first-best welfare. I build reproducible pipelines on AWS and HPC/SLURM in Python, R, and C++, and write tests to keep the code clear and correct. I also apply merger simulation, market entry and exit counterfactuals, and firm-conduct identification.

My CV lists an expected graduation of June 2027, but if an offer is available I can finish earlier and join as early as November or December for a Ph.D.\ Associate or Analyst role. I am also open to an Intern position; for that, timing may need to wait until next year, or I could start remotely if helpful. I have attached my CV and transcript, and included my job market paper as a writing sample and my GitHub (\href{https://github.com/cliohsu1997}{github.com/cliohsu1997}). Thank you for your consideration.

\vspace{0.7em}
Sincerely,\\[1.0em]
Ke-Cheng Hsu (Clio)

\end{document}
